\section{Análise dos Stakeholders}\label{sec:stakeholders}
A Tabela~\ref{tab:stakeholders} organiza seis papéis. As funcionalidades,
contribuições e testes sustentam sua identificação; necessidades e conflitos
são inferências analíticas, sem entrevistas. Não se afirma a existência de
seis equipes separadas na organização do projeto.

\begin{table*}[t]
\caption{Stakeholders: evidência do papel e necessidades inferidas.}
\label{tab:stakeholders}
\centering
\begin{tabularx}{\textwidth}{@{}lXX@{}}
\toprule
Parte interessada & Evidência observada & Interesse e necessidade inferidos \\
\midrule
Usuários & Leitura, tópicos e recursos salvos & Conteúdo pertinente, navegação clara e acesso local \\
Desenvolvedores & Módulos, código e PRs & Compreender limites, implementar mudanças e testar isoladamente \\
Mantenedores & Revisões, documentação e automação & Conservar coerência, estabilidade e builds viáveis \\
Produtores de conteúdo & Série NiA referenciada pelo aplicativo & Conteúdo correto e atualizado; processo editorial não auditado \\
Qualidade & Testes JVM, instrumentados e screenshots & Reproduzir falhas e verificar estados e persistência \\
Plataforma Android & APIs Jetpack e objetivo de referência & Evoluir práticas e manter exemplos compreensíveis \\
\bottomrule
\end{tabularx}
\end{table*}

\begin{table}[htb]
\caption{Conflitos possíveis e encaminhamentos propostos.}
\label{tab:conflitos}
\centering
\begin{tabularx}{\columnwidth}{@{}XX@{}}
\toprule
Interesses em tensão & Encaminhamento \\
\midrule
Atualização do conteúdo versus rede e bateria & Leitura local e trabalho de sincronização agendado \\
Novas APIs versus estabilidade e compreensão & Versões fixadas, revisão e testes de regressão \\
Isolamento das features versus custo de configuração & Módulos proporcionais ao escopo e regras de dependência \\
\bottomrule
\end{tabularx}
\end{table}

Esses conflitos não são incidentes comprovados. São hipóteses úteis para
justificar critérios de avaliação e explicitar os custos das decisões.
